% Part: normal-modal-logic
% Chapter: frame-correspondence
% Section: first-order-definability

\documentclass[../../../include/open-logic-section]{subfiles}

\begin{document}

\olfileid{nml}{frd}{fol}

\olsection{پہلے درجے دی زبان وچ تعریف}

اسیں ویکھیا اے کہ فریماں دے رسائی دے تعلقاں دیاں کئی خاصیتاں امکانی !!{formula}s نال متعین ہو سکدیاں نیں۔ مثال لئی فریماں دا تقابل !!{formula}~\Ax{B}، یعنی $p \lif \Box\Diamond
p$، نال متعین ہو سکدا اے۔ ہُن تک جیہڑیاں شرطاں ویکھیاں نیں اوہ ساریاں پہلے درجے دے !!{formula}s نال ایسی !!a{language} وچ دَسیاں جا سکدیاں نیں جیہدے وچ صرف اک دو رکنی !!{predicate} ہووے۔ مثال لئی تقابل
$\lforall[x][\lforall[y][(\Atom{Q}{x,y} \lif \Atom{Q}{y, x})]]$ نال متعین ہوندا اے، ایس معنی وچ کہ پہلے درجے دی !!{structure}~$\Struct{M}$، جیہدے وچ $\Domain{M} =
W$ تے $\Assign{Q}{M} = R$ اے، اُپرلا فارمولا اُدوں تے صرف اُدوں پورا کردی اے جدوں $R$ متقابل ہووے۔ ایس توں ہِیٹھلی تعریف دا خیال آؤندا اے۔

\begin{defn}
  فریماں دی جماعت~$\mClass{F}$ \emph{پہلے درجے دی زبان وچ متعین ہو سکدی اے} جے کوئی ایسا !!a{sentence}~$!A$ ہووے، پہلے درجے دی اوس زبان وچ جیہدے وچ اک دو رکنی !!{predicate}~$Q$ ہووے، تے ایہ شرط پوری ہوندی ہووے کہ $\mModel{F} =
  \tuple{W, R} \in \mClass{F}$ اُدوں تے صرف اُدوں جدوں $\Sat{M}{!A}$ ہووے، پہلے درجے دی اوس !!{structure}~$\Struct{M}$ وچ جیہدے لئی $\Domain{M} = W$ تے $\Assign{Q}{M}
  = R$۔
\end{defn}

پتا لگدا اے کہ ہُن تک ویکھیاں خاصیتاں تے اوہناں نوں متعین کرن والے امکانی !!{formula}s خاص معاملے نیں۔ کجھ !!{formula} ایسے نیں جیہڑے فریماں دی ایسی جماعت متعین کردے نیں جیہڑی پہلے درجے دی زبان وچ متعین نہیں ہو سکدی۔ تے فریماں دیاں کجھ اوہ جماعتاں، جیہڑیاں پہلے درجے دی زبان وچ متعین ہو سکدیاں نیں، کسے امکانی فارمولے نال متعین نہیں ہو سکدیاں۔

پہلی گل دی ردّی مثال لوب دا فارمولا اے۔
\begin{equation}
\Box(\Box p \lif p) \lif \Box p. \tag{\Ax{W}}
\end{equation}
\Ax{W} اوہناں فریماں دی جماعت متعین کردا اے جیہڑے متعدی تے اُلٹے رخ توں پکّی بنیاد والے نیں۔ تعلق پکّی بنیاد والا اے جے کوئی لامتناہی سلسلہ $w_1$، $w_2$، \dots{} ایسا نہ ہووے جیہدے لئی $Rw_2w_1$، $Rw_3w_2$، \dots۔ مثلاً تعلق $<$،~$\Nat$ اُتے پکّی بنیاد والا اے، جدکہ تعلق $<$،~$\Int$ اُتے پکّی بنیاد والا نہیں۔ تعلق اُلٹے رخ توں پکّی بنیاد والا اُدوں تے صرف اُدوں اے جدوں اوہدا اُلٹا تعلق پکّی بنیاد والا ہووے۔ سو اُلٹے رخ توں پکّی بنیاد والے تعلق اوہ نیں جیہناں وچ کوئی لامتناہی سلسلہ $w_1$، $w_2$، \dots{} ایسا نہیں ہوندا جیہدے لئی $Rw_1w_2$، $Rw_2w_3$، \dots۔

پر پہلے درجے دا کوئی !!{formula} نہیں جیہڑا متعدی تے اُلٹے رخ توں پکّی بنیاد والے تعلقاں نوں متعین کرے۔ منّ لو~$\Sat{M}{!F}$ اُدوں تے صرف اُدوں جدوں $R =
\Assign{Q}{M}$ متعدی تے اُلٹے رخ توں پکّی بنیاد والا ہووے۔ منّ لو $!A_n$ ایہ !!{formula} اے۔
\[
(\Atom{Q}{a_1,a_2} \land \dots \land
    \Atom{Q}{a_{n-1},a_{n}})
\]
ہُن ایہ !!{formula}s دا سیٹ ویکھو۔
\[
\Gamma = \{!F, !A_1, !A_2, \dots\}.
\]
$\Gamma$ دا ہر محدود ذیلی سیٹ پورا ہون جوگا اے۔ سبھ توں وڈا $k$ لے لو جیہدے لئی $!A_k$ ذیلی سیٹ وچ اے، تے $\Domain{M_k} = \{1, \dots, k\}$،
$\Assign{a_i}{M_k} = i$ تے $\Assign{Q}{M_k} = <$ رکھو۔ کیوں جے $<$، $\{1,
\dots, k\}$ اُتے متعدی تے اُلٹے رخ توں پکّی بنیاد والا اے، ایس لئی
$\Sat{M_k}{!F}$۔ بناوٹ مطابق $\Sat{M_k}{!A_i}$ سارے $i \le k$ لئی۔ پہلے درجے دی منطق دے محدود گواہی دے قضیے مطابق $\Gamma$ کسے !!{structure}~$\Struct{M}$ وچ پورا ہون جوگا اے۔ مفروضے مطابق، کیوں جے
$\Sat{M}{!F}$، ایس لئی تعلق $\Assign{Q}{M}$ اُلٹے رخ توں پکّی بنیاد والا اے۔ پر صاف اے کہ $\Assign{a_1}{M}$، $\Assign{a_2}{M}$،
\dots{} اوہو جیہا لامتناہی سلسلہ بناؤندے نیں جیہڑا اُلٹے رخ توں پکّی بنیاد ہون دی شرط منع کردی اے۔

دوجے دعوے دی ردّی مثال ہمہ گیری دی خاصیت اے۔ ہر $u$ تے $v$ لئی $Ruv$۔ ہمہ گیر فریم پہلے درجے دے فارمولے
$\lforall[x][\lforall[y][\Atom{Q}{x,y}]]$ نال متعین ہوندے نیں۔ پر کوئی امکانی فارمولا ایہو جیہا نہیں جیہڑا سارے ہمہ گیر فریماں وچ، تے صرف اوہناں وچ، معتبر ہووے۔ ایہ اک نتیجے توں نکلدا اے جیہڑا آپ وی دلچسپ اے۔ ہمہ گیر فریماں وچ معتبر فارمولے عین اوہی نیں جیہڑے انعکاسی، متقابل تے متعدی فریماں وچ معتبر نیں۔ ایسے انعکاسی، متقابل تے متعدی فریم نیں جیہڑے ہمہ گیر نہیں۔ ایس لئی ہر اوہ !!{formula} جیہڑا سارے ہمہ گیر فریماں وچ معتبر اے، کجھ غیر ہمہ گیر فریماں وچ وی معتبر اے۔

\end{document}
